\section{Shell counts and smearing moments}\label{sec:moments}

All distances in this section are in lattice units.  A \emph{smearing of
radius $a\geq1$ and weight bound $A$} is a linear statistic
\begin{equation}\label{moments:weights}
 Y(H)=\sum_u w_uH(u),\qquad
 \supp w\subset B(z_0,a),\qquad
 \sum_u|w_u|\leq A,\qquad \sup_u|w_u|\leq Aa^{-2}.
\end{equation}
The weights are real and may have either sign.  If $\sum_uw_u=0$, we
call the smearing \emph{zero-total}; its value is independent of the
additive constant in the height lift.  Constants throughout may depend
on the fixed parameter $x$.

The bounded height problems below use the convention of
Lemma~\ref{lem:height-order}: a connected fixed-trace lattice skeleton
bounds the free sites, and every incident interaction up to that trace
is retained.  This includes the interior of a stopped contour after its
inner face layer has been fixed.  The skeleton need not be convex or
have a lower bound on the widths of its fjords.

We will prove two bounds for these smearings, for every finite $p\geq1$.
Suppose the radius $a$ is at most a fixed small fraction of the distance
$n$ from the smearing center to the trace.  If that trace has magnitude
at most $K$, the $L^p$ norm is controlled by
$1+K+\sqrt{\log(n/a)}$.  For zero-total smearings in the free plane
law, the norms are bounded independently of $a$.  The common starting point is a bound on the
number of enclosing loops in each shell.  It gives a logarithmic
covariance estimate and hence positive mass in a central interval;
rounded midpoints then turn that central mass into all finite moments.

\subsection{A shell estimate without an upper crossing input}

\begin{lemma}[Geometric tails for shell counts]\label{lem:shell-moments}
In a finite constant-boundary height problem, let $N_d(z)$ be the number
of level loops surrounding a site $z$ whose Euclidean closest approach
to $z$ belongs to $[d,2d)$, where $d\geq1$.  There are $c,C>0$ such that
\begin{equation}\label{moments:shell-tail}
 \Prob\{N_d(z)\geq t\}\leq Ce^{-ct}\qquad(t\geq0).
\end{equation}
In particular, for each finite $p\geq1$, $\E N_d(z)^p\leq C_p$.
These constants are independent of the shape, $z$, and $d$, and the
same assertions hold in any bounded face of a connected flat skeleton.
\end{lemma}

\begin{proof}
A constant shift reduces to boundary height zero.  Conditional on their
unoriented geometry, the level loops have independent fair height
increments.  Expose enclosing loops from outside inward, using the
stopping rule in Lemma~\ref{lem:height-order}.  If the first loop counted
by $N_d(z)$ exists, its unexplored interior has the empty-exterior loop
law with a constant inner trace.  Its distance from $z$ is at most
$2d+O(1)$.  Every further counted loop surrounds $B(z,d/2)$, for all
sufficiently large $d$.  The factor $1/2$ provides clearance from all
rounding and incident-face conventions.

We first prove a uniform obstruction in any such remaining domain.
Suppose an exterior search finds a good open-plus $B$ contour $C$ strictly
inside its boundary, surrounding $B(z,d/2)$.  Choose the outermost such
contour by the authorized-edge invasion following
Lemma~\ref{discrete:online}.  The search exposes only exterior
history and the prescribed open-plus cut data.  In particular it reveals
no strict interior cell and leaves the exact cut law of
Lemma~\ref{lem:cell-cut} inside $C$.  Moreover
\begin{equation}\label{moments:cut-radius}
 d/2\leq\dist(z,C)\leq2d+O(1).
\end{equation}
There is a number $\kappa>0$, independent of $C$, such that under this
interior law a $B$-plus path joins its prescribed inner layer to
$B(z,d/2)$ with probability at least $\kappa$.

Here are geometric and conditional details of this last assertion.
Fix a sufficiently small constant $\epsilon>0$, and choose a closest
point $p$ of $C$ to $z$.  Realize the same interior law in a new free
ambient triangulation containing $B(z,10d)$, conditioned only on the
same open-plus cells and first-spin sites of $C$.  Exact cutting
justifies this replacement even for an arbitrarily nonconvex $C$.
In this ambient graph request a good plus circuit in an annulus about
$p$ with radii $\epsilon d$ and $2\epsilon d$.  Since $C$ visits the
inner ball and surrounds $B(z,d/2)$, it must exit the outer ball and
intersect this circuit.  Next request a chain of annular plus circuits
of these same radii, from $p$ to an annulus wholly inside $B(z,d/2)$.
The centers may follow a straight segment in the ambient plane.  Choose
successive spacings strictly between $\epsilon d$ and $2\epsilon d$,
with fixed slack.  The inner balls of consecutive annuli overlap and
each protrudes outside the other's outer ball.  Consequently their
surrounding circuits intersect: disjoint circuits would have either
disjoint interiors, contradicting the overlap, or nested interiors,
contradicting that protrusion.  Equation~\eqref{moments:cut-radius}
bounds the number of requests by a constant depending only on
$\epsilon$.  Bounded lattice displacements are absorbed by the slack.

Lemma~\ref{lem:buffer-circuits} and positive association give a fixed
positive probability for their intersection in the free ambient law.
All the data imposed on $C$ are favorable for this increasing event,
so Proposition~\ref{prop:history-order} gives the same lower bound after
conditioning on $C$.  The connected circuit chain contains a path
from $C$ to the terminal circuit inside the disk.  Take its last
intersection with $C$ before that endpoint.  The remaining path stays
inside $C$, since otherwise planar separation would force another
intersection.  Its adjacent plus face chain therefore joins the
prescribed inner layer to the disk.  This proves the asserted lower
bound $\kappa$.  In particular this connection precludes a deeper
$B$ wall surrounding the disk.  The ambient chain was allowed to leave
$C$; no interior corridor of prescribed width was assumed.

To turn this obstruction into a bound on the number of loops, let $N$
be the number of further qualifying enclosing loops in the remaining
flat problem.  Closest approach is nonincreasing under nesting, so the
qualifying loops are consecutive in the enclosing nesting order until
their closest approach falls below $d$.  Conditional on all unoriented
geometry, write the height residues along them as a nearest-neighbor
walk on $\Z/4\Z$, with independent fair signs.  From every starting
residue there is a prescribed block of four signs that contains a
$W$ wall with unchanged $B=+$, followed inward by a $B$ wall.  For
example, starting from residues $0,1,2,3$, respectively, the initial
strings $++$, $--$, $---$, $+++$ have that property; complete each to
four signs arbitrarily.  Every chosen four-sign block has probability
$1/16$.  Hence, on $\{N\geq k\}$, the conditional probability of
seeing this ordered pair of walls among the first $k$ loops is at least
\[
 1-\varepsilon_k,
 \qquad \varepsilon_k=(15/16)^{\lfloor k/4\rfloor}.
\]
The statement follows by conditioning successively on the starting
residue of each block; independence between block-success events is
not required.

Call this ordered-wall event $\mathcal A$.  A $W$ wall with $B=+$ is
a good open-plus contour, because every cell at which $W$ is
nonconstant is necessarily open.  Thus $\mathcal A$ implies existence
of the exterior-searched outermost plus cut above and of a $B$ wall
strictly inside it surrounding the disk.  The preceding obstruction,
conditioned on every possible stopped cut, yields
\[
 (1-\varepsilon_k)\Prob(N\geq k)
 \leq\Prob(\mathcal A)\leq1-\kappa.
\]
Choose a fixed integer $k$ with $\varepsilon_k\leq\kappa/2$ and put
$\rho=(1-\kappa)/(1-\kappa/2)<1$.  The probability of continuing for
$k$ further qualifying loops is at most $\rho$.

After exposing these $k$ loops, the last one's interior is again an
empty-exterior flat problem of exactly the same type, with closest
boundary distance at most $2d+O(1)$.  Restarting conditionally gives
\[
 \Prob\{N_d(z)\geq1+mk\}\leq\rho^m\qquad(m\geq0).
\]
This is a stopping recursion; it asserts no independence of successive
geometric shells.  It proves~\eqref{moments:shell-tail} for sufficiently
large $d$.  For bounded $d$, each counted loop visits a fixed-size ball
and the loops have disjoint honeycomb vertices, giving a deterministic
bound on their number.  Finally, integration of the exponential tail
gives every stated moment.
\end{proof}

\begin{lemma}[Common loops and the logarithmic covariance bound]
\label{moments:covariance}
Let the flat trace be zero, let $n$ be the distance from $z_0$ to its
bounding skeleton, and let $z,w\in B(z_0,a)$, where
$1\leq a\leq c_0n$ for a sufficiently small absolute $c_0$.  With
$b=\max(1,|z-w|)$,
\begin{equation}\label{moments:covariance-bound}
 0\leq\E[H(z)H(w)]
 \leq C\left(1+\log\frac{n}{b}\right).
\end{equation}
Every smearing~\eqref{moments:weights} in this ball therefore satisfies
\begin{equation}\label{moments:second-moment}
 \E Y^2\leq C A^2\left(1+\log\frac{n}{a}\right).
\end{equation}
\end{lemma}

\begin{proof}
The independent centered loop orientations give the exact identity
\begin{equation}\label{moments:common-loops}
 \E[H(z)H(w)]
 =\E\#\{\text{level loops surrounding both $z$ and $w$}\}.
\end{equation}
Fix a dyadic closest-approach shell $[d,2d)$ about $z$ with
$d\leq b/100$.  A loop surrounding both sites and entering $B(z,2d)$
must also leave a ball about $z$ of radius comparable to $b$; otherwise
it could not enclose $w$.  It supplies a wall arm across the intervening
annuli.  Along the whole loop exactly one of $B,W$ changes.  A good
circuit of that color cannot be crossed by its wall, since the circuit's
open cells have constant first spin.

Use the wall color as the first spin.  The zero-height skeleton fixes
that first spin and imposes constancy of the second spin along a
connected skeleton.  Represent the latter constancy by $E_0$ edges.
Its actual constant label costs exactly $1/2$, independent of the
first spins and statuses, and hence does not change their marginal.
All arm annuli lie away from the skeleton, since $b\leq2a+1$ and
$c_0$ is small.  The logarithmic circuit stack from
Lemma~\ref{lem:buffer-circuits}, applied for each of the two possible
wall colors, consequently gives
\begin{equation}\label{moments:arm}
 \Prob\{\text{such a common loop exists in the shell}\}
 \leq C(d/b)^\alpha.
\end{equation}
If $M_d$ counts common loops in this shell, then
$M_d\leq N_d(z)$ and $M_d=0$ outside the event in~\eqref{moments:arm}.
Cauchy--Schwarz and Lemma~\ref{lem:shell-moments} imply
\[
 \E M_d\leq (\E N_d(z)^2)^{1/2}
       \Prob(M_d>0)^{1/2}
 \leq C(d/b)^{\alpha/2}.
\]
Loops at distance below one give a deterministic bounded contribution
by disjointness.  The sum over the small dyadic scales is bounded.
Each of the
$O(1+\log(n/b))$ remaining shells contributes at most a constant.
There are no shells beyond scale $O(n)$: a nearest skeleton site is
outside every enclosing loop, so the segment from $z$ to that site
meets the loop within distance $n+a+O(1)$.  These facts and
\eqref{moments:common-loops} prove~\eqref{moments:covariance-bound}.

For completeness the average of the singularity is uniformly bounded,
including its diagonal.  Triangular-lattice balls have at most
$C(1+r)^2$ sites.  For fixed $z$, the layer-cake formula gives, when
$a>1$,
\begin{align*}
 &\sum_w |w_w|\log^+\frac{a}{\max(1,|z-w|)}\\
 &\qquad=\int_0^{\log a}
      \sum_{\max(1,|z-w|)<ae^{-t}}|w_w|\,\dd t
 \leq C A a^{-2}\int_0^{\log a}a^2e^{-2t}\,\dd t
 \leq C A.
\end{align*}
For $a=1$ the sum vanishes.  Sum
\eqref{moments:covariance-bound} against $|w_zw_w|$, separate the
constant term $\log(n/a)$, and use this bound and
$\sum|w_z|\leq A$.  This proves~\eqref{moments:second-moment}
for arbitrary signed weights.
\end{proof}

\subsection{Weighted midpoint amplification}

We use the discrete midpoint inequality of Klartag--Lehec
\cite{KlartagLehec2019} in the nonnegative-function formulation of
Gozlan--Roberto--Samson--Tetali \cite[Theorem~3]{GRST2019}.  Its
one-dimensional finite-support statement follows directly from their
entropy inequality \cite[Theorem~8, equation~(9)]{GRST2019}, including
zero weights.  We give that deduction, its coordinate induction, and
all additional hypotheses used here.

\begin{theorem}[Discrete midpoint inequality]\label{thm:midpoint}
Let $f,g,u,v:\Z^m\longrightarrow[0,\infty)$ have finite support, and
write $m_-(s,t)=\lfloor(s+t)/2\rfloor$ and
$m_+(s,t)=\lceil(s+t)/2\rceil$, with coordinatewise rounding.
If, for every $s,t\in\Z^m$,
\[
 f(s)g(t)\leq u(m_-(s,t))v(m_+(s,t)),
\]
then
\[
 \left(\sum f\right)\left(\sum g\right)
 \leq\left(\sum u\right)\left(\sum v\right).
\]
\end{theorem}

\begin{proof}
First let $m=1$.  Write $A=\sum f$ and $B=\sum g$.
If either is zero there is nothing to prove.  Otherwise let
$\nu_0=f/A$ and $\nu_1=g/B$, and couple them monotonically: take
their quantile functions at the same uniform point of $(0,1)$.
Let $\pi$ be this coupling, and let $\nu_-,\nu_+$ be its pushforwards
by the floor and ceiling midpoint maps.  The entropy theorem just
cited states
\[
 \mathcal H(\nu_-)+\mathcal H(\nu_+)
 \le\mathcal H(\nu_0)+\mathcal H(\nu_1),
 \qquad \mathcal H(\nu)=\sum_z\nu(z)\log\nu(z),
\]
for finitely supported probability measures on $\Z$, with zero terms
omitted.  Every pair charged by $\pi$ has $f(s)g(t)>0$; the
pointwise hypothesis therefore makes both corresponding output weights
positive.  All logarithms below are thus evaluated on finite positive
supports.  Integrating the logarithm of the pointwise hypothesis and
then using the entropy inequality gives
\begin{align*}
 \log A+\log B
 &\le\sum_z\nu_-(z)\log u(z)+\sum_z\nu_+(z)\log v(z)
            -\mathcal H(\nu_0)-\mathcal H(\nu_1)\\
 &\le\left(\sum_z\nu_-(z)\log u(z)-\mathcal H(\nu_-)\right)
     +\left(\sum_z\nu_+(z)\log v(z)-\mathcal H(\nu_+)\right)\\
 &\le\log\sum_z u(z)+\log\sum_z v(z).
\end{align*}
The last inequality is Jensen's inequality applied to
$\sum\nu\log(a/\nu)$ on the positive support of each probability
measure; adding any unused nonnegative weights only enlarges the sum.
Exponentiating proves the scalar statement without modifying zero
values or assuming that a positive approximation preserves its premise.

For the induction step write $s=(s',a)$ and $t=(t',b)$.  Fix $s',t'$
and apply the one-dimensional theorem in $a,b$.  For the partial sums
$F(s')=\sum_a f(s',a)$, and likewise $G,U,V$, its conclusion is
\[
 F(s')G(t')\leq
 U(m_-(s',t'))V(m_+(s',t')).
\]
The induction hypothesis in the other $m-1$ coordinates proves the
claim.  No factorization of the functions is used.  A finite array is
extended by zero to $\Z^m$; the premise must still be checked on every
pair of its nonzero input supports.
\end{proof}

\begin{lemma}[Admissibility and energy at rounded midpoints]
\label{moments:weighted-midpoints}
Let $\mathcal S$ be the integer $1$-Lipschitz arrays on a finite
connected graph, with fixed integer heights on a nonempty set of
vertices.  Set
\[
 \mathcal W(h)=\ind_{\{h\in\mathcal S\}}
       x^{\sum_{\{i,j\}}|h_i-h_j|}.
\]
Then $\mathcal S$ is finite, both rounded midpoints of $s,t\in\mathcal S$
belong to $\mathcal S$, and
\begin{equation}\label{moments:weight-midpoint}
 \mathcal W(m_-(s,t))\mathcal W(m_+(s,t))\geq \mathcal W(s)\mathcal W(t).
\end{equation}
Moreover a linear statistic with $\sum|w_i|\leq A$ obeys
\begin{equation}\label{moments:rounding}
 \left|Y(m_\pm(s,t))-\frac{Y(s)+Y(t)}2\right|\leq A/2.
\end{equation}
The same energy inequality holds with additional support constraints
whenever both outputs satisfy those constraints.
\end{lemma}

\begin{proof}
One fixed height bounds every coordinate by its graph distance from
that vertex, so the array is finite.  For an edge $ij$ put
$a=s_i-s_j$ and $b=t_i-t_j$.  If $c,d$ are its differences in the two
rounded arrays, and $\varepsilon_i$ is the parity of $s_i+t_i$, then
\[
 c=\frac{a+b-\varepsilon_i+\varepsilon_j}{2},\qquad
 d=\frac{a+b+\varepsilon_i-\varepsilon_j}{2}.
\]
Thus $c+d=a+b$, $|c-d|\leq1$, and $c,d$ are its two nearest integer
halves.  In particular $c,d\in\{-1,0,1\}$ and
\[
 |c|+|d|=|a+b|\leq|a|+|b|.
\]
Each fixed integer coordinate is preserved exactly.  Summing over edges
and using $x\leq1$ proves~\eqref{moments:weight-midpoint}, including
$x=1$.  Each rounded coordinate differs from its average by at most
$1/2$, which proves~\eqref{moments:rounding} even for signed weights.
\end{proof}

\begin{lemma}[Central mass in finitely many classes]
\label{lem:midpoint-amplification}
Let $\mu(h)=\mathcal W(h)/Z$ be a finite probability measure on integer arrays.
Partition its support into at most $J$ classes.  Suppose that two arrays
from the same class have both rounded midpoints in the full support,
with the weight inequality~\eqref{moments:weight-midpoint}.  Let $Y$ be
a linear statistic satisfying~\eqref{moments:rounding}.  If $M>0$,
$M\geq A/2$, and, for every class $\mathcal S_i$,
\begin{equation}\label{moments:class-central}
 \mu\bigl(\mathcal S_i\cap\{|Y|\leq M\}\bigr)\geq c>0,
\end{equation}
then there are $a_*>0$ and $C_p<\infty$, depending only on $J,c,p$, such
that
\begin{equation}\label{moments:amplification}
 \mu\{|Y|>tM\}\leq4e^{-a_*t}\quad(t\geq0),
 \qquad \|Y\|_{L^p(\mu)}\leq C_pM\quad(1\leq p<\infty).
\end{equation}
Outputs may belong to any class.  In particular no stability of an
individual class under midpoint rounding is required.
\end{lemma}

\begin{proof}
For integers $j\geq3$ define
\[
 I_j=[(2^j-4)M,(2^{j+1}+4)M],\qquad p_j=\mu(Y\in I_j).
\]
For $j\geq4$ choose a class carrying at least $p_j/J$ of the mass of
$I_j$, and pair its tail arrays with its central arrays from
\eqref{moments:class-central}.  Every output has its statistic in
\[
 [(2^{j-1}-7/2)M,(2^j+7/2)M]\subset I_{j-1};
\]
indeed the central average contributes at most $M/2$ and rounding at
most $A/2\leq M$.  Apply Theorem~\ref{thm:midpoint} to $\mathcal W$ restricted
to the chosen tail and central sets, and to two copies of $\mathcal W$ restricted
to $\{Y\in I_{j-1}\}$.  Dividing by $Z^2$ gives
\begin{equation}\label{moments:recursion}
 p_j\leq (J/c)p_{j-1}^2.
\end{equation}

Put $K_*=J/c\geq1$, $N=\lceil4K_*\rceil$, and $L=3+N$.
For $j\geq4$ these intervals have overlap at most two: the upper
endpoint of $I_j$ is strictly below the lower endpoint of $I_{j+2}$.
Thus $\sum_{j\geq4}p_j\leq2$, and among $j=4,\ldots,L$ there is
$j_0$ with $K_*p_{j_0}\leq1/2$.  Iteration of
\eqref{moments:recursion} yields
\[
 K_*p_{j_0+r}\leq2^{-2^r}\quad(r\geq0),
 \qquad
 p_j\leq K_*^{-1}
 \exp\{- (\log2)2^{j-L}\}\quad(j\geq L).
\]
Repeat the argument for $-Y$, with the same $L$.  If $t\geq2^L$ and
$k=\lfloor\log_2t\rfloor$, the cores
$[2^jM,2^{j+1}M]\subset I_j$ for $j\geq k$ cover the positive tail,
and their negative counterparts cover the negative tail.  Writing
$q_0=\exp\{-(\log2)2^{k-L}\}\leq1/2$, we obtain
\[
 \mu(|Y|>tM)
 \leq\frac2{K_*}\sum_{r\geq0}q_0^{2^r}
 \leq4q_0
 \leq4\exp\{- (\log2)2^{-L-1}t\}.
\]
For $0\leq t<2^L$ the last bound exceeds one.  This proves the tail
claim with $a_*=(\log2)2^{-L-1}$.  Finally
$\E|Y|^p=\int_0^\infty pt^{p-1}\mu(|Y|>t)\,\dd t$ proves the moment
claim.  This proof explicitly uses the unconditional central mass in
each class; a central bound for their union would not suffice.
\end{proof}

\subsection{Bounded traces and plane increments}

\begin{proposition}[Smearing with bounded boundary values]
\label{prop:smearing}
In a bounded height problem as specified above, suppose its fixed trace
has absolute value at most $K$.  Let $n$ be the distance from the
smearing center to the bounding sites.  For
$1\leq a\leq c_0n$ and every finite $p\geq1$,
\begin{equation}\label{moments:smearing-bound}
 \left\|\sum_uw_uH(u)\right\|_p
 \leq C_{p,A}\left(1+K+\sqrt{\log(n/a)}\right).
\end{equation}
The bound is uniform over the shape of the trace and applies in
particular to the interior of a good cut after its complete inner
height layer has been specified.
\end{proposition}

\begin{proof}
For zero trace, Lemma~\ref{moments:covariance} and Chebyshev's inequality
supply central probability at least $1/2$ with
$M=C_A(1+\sqrt{\log(n/a)})$, increased if necessary to exceed $A/2$.
Lemma~\ref{moments:weighted-midpoints} verifies all midpoint hypotheses
for the actual weighted height measure.  Apply
Lemma~\ref{lem:midpoint-amplification} with one class to prove
\eqref{moments:smearing-bound} for zero trace.

For general trace, put $K_0=\lceil K\rceil$.  Height order from
Lemma~\ref{lem:height-order} traps the law between the flat laws with
traces $-K_0$ and $K_0$.  For nonnegative weights, the corresponding
smearing laws are the zero-trace smearing shifted by
$\pm K_0\sum_uw_u$.  Both tails are therefore bounded by the
zero-trace tails with an additive displacement at most $AK_0$.
This gives~\eqref{moments:smearing-bound}.  Split arbitrary signed
weights into their positive and negative parts; each satisfies
\eqref{moments:weights}, and the triangle inequality in $L^p$ finishes
the proof.  On a cut, fix the entire actual inner trace first, so its
interior is an ordinary weighted height problem by
Lemma~\ref{lem:cell-cut}.  The same argument applies conditionally.
\end{proof}

\begin{proposition}[Uniform plane moments for zero-total smearings]
\label{prop:plane-moments}
For the infinite-volume increment law furnished by
Proposition~\ref{prop:local-mixing}, every zero-total
smearing~\eqref{moments:weights} satisfies
\begin{equation}\label{moments:plane-bound}
 \left\|\sum_uw_uH(u)\right\|_p\leq C_{p,A}
 \qquad(1\leq p<\infty).
\end{equation}
The constants are uniform in $a\geq1$ and the center.  The same estimate
holds for the increment law in a finite free simply connected
triangulation containing $B(z_0,C_0a)$, for a sufficiently large fixed
$C_0$.
\end{proposition}

\begin{proof}
Work first with the free spin pair in the finite ambient graph.
Choose a fixed large constant $L_0$ with $1/L_0<c_0/4$, and search
from outside for a good $B$ circuit in an annulus with radii
$L_0a$ and $2L_0a$.  Choose $C_0$ large enough that all search buffers
fit in the ambient graph.  Lemma~\ref{lem:buffer-circuits} gives this
stopped-cut event probability at least $c_1>0$.  Its trace is contained
in two consecutive integers.  Conditional also on the complete inner
trace, subtract the lower of these integers.  The interior is now an
ordinary height problem with trace in $\{0,1\}$, and its distance from
$z_0$ is between fixed multiples of $L_0a$.  Proposition~\ref{prop:smearing}
therefore gives a uniform conditional second-moment bound $C_A$.
The random subtraction changes no zero-total smearing.  Chebyshev's
inequality and then averaging over the search show, for some fixed
$M\geq A/2$,
\begin{equation}\label{moments:plane-central}
 \Prob(|Y|\leq M)\geq c_1/2.
\end{equation}
If the finitely many smallest lattice scales require enlarging the
annulus to accommodate rounding, use a fixed minimal annulus radius;
its ratio to $a\geq1$ remains bounded and the same constants suffice.

To amplify this central bound, use the correct finite height measure.
For any chosen anchor $o$, the four cyclic rotations of the Gray states
preserve every increment and every interaction weight, and act
transitively on the pair state at $o$.  Consequently the unconditioned
free pair has exactly the same increment law as its version with state
$(+,+)$ at $o$.  By Lemma~\ref{lem:spin-lift}, the latter is the measure
on integer $1$-Lipschitz arrays with $H(o)=0$ and density proportional
to $x^{\sum_{ij}|H(i)-H(j)|}$.  Thus $Y$ in this anchored measure has
the central bound~\eqref{moments:plane-central}.  This argument obtains
the cut central bound before fixing the anchor, and hence introduces
no extra anchor constraint into the stopped interior law.

The anchored arrays form one midpoint-admissible class by
Lemma~\ref{moments:weighted-midpoints}.  Applying
Lemma~\ref{lem:midpoint-amplification} proves~\eqref{moments:plane-bound}
in finite volume.  At fixed lattice weights, enlarge that volume.
The increment convergence from Proposition~\ref{prop:local-mixing}
and the uniform bounds pass the estimate to the plane; higher moments
also give convergence of every fixed finite moment.  All constants
above are independent of the location and radius of the smearing.
\end{proof}

\begin{lemma}[Nondegeneracy between two disks]\label{lem:nondegeneracy}
Let $U_1,U_2$ be two fixed disjoint round disks with disjoint closures,
and let $f_i\geq0$ be smooth functions supported compactly in $U_i$,
with $\int f_i=1$.  Form their exact face-area averages on the dilated
lattice geometry of scale $R$.  If $Y_R$ is the first average minus
the second, there is $c>0$ such that
\begin{equation}\label{moments:nondegeneracy}
 \E Y_R^2\geq c
\end{equation}
for all sufficiently large $R$, in the plane and in sufficiently large
finite free volumes.  The bound is uniform under translations of the
fixed arrangement.
\end{lemma}

\begin{proof}
Choose two disjoint nested annular collars around $\supp f_1$ inside
$U_1$, avoiding $\supp f_2$.  Prescribe an outer good $B$-plus cut and
an inner good $B$-minus cut.  Successive applications of
Lemma~\ref{lem:buffer-circuits} in these separated buffers give their
joint probability at least $c_2>0$, uniformly in $R$ large enough for
rounding.  Make the choices by the stopped searches following
Lemma~\ref{discrete:online}; the event and the selected cuts
depend on first spins and cell statuses, and do not prescribe any
second-spin labels.

By Lemma~\ref{lem:collar-switch}, conditional on those variables there
is a second-spin component in the intervening collar whose fair label
flip changes every height in the inner protected region by $4$ in
absolute value relative to the exterior of the outer cut.  Choose the
component measurably using a fixed ordering, and condition on all its
other labels.  The cuts, their event, and all weights of $Y_R$ are
unchanged by this flip.  Exact unit mass in the first disk and zero
weight in the collar mean that the two values of $Y_R$ differ by $4$.
A fair pair of real numbers at distance $4$ has variance $4$ and
second moment at least $4$.  Averaging over the cut event yields
$\E Y_R^2\geq4c_2$ in finite volume.  The fourth-moment bound in
Proposition~\ref{prop:plane-moments} allows passage of these second
moments to the plane.  This proves the claim.
\end{proof}

\subsection{Exact area interpolation}

Let $Q_{\delta,u}$ be the hexagonal face centered at $u\in\delta\T$.
For a bounded compactly supported test $f$, use precisely
\begin{equation}\label{moments:area-weights}
 w_{\delta,f}(u)=\int_{Q_{\delta,u}}f(z)\,\dd z,
 \qquad
 \langle h_\delta,f\rangle
 =\sum_u w_{\delta,f}(u)h_\delta(u).
\end{equation}
This identity follows from the piecewise constant extension in the
problem, including its zero extension outside the discrete domain.
For tests supported inside that domain the tessellation gives the
exact, rather than asymptotic, identities
\[
 \sum_u w_{\delta,f}(u)=\int f,
 \qquad \sum_u|w_{\delta,f}(u)|\leq\|f\|_{L^1},
 \qquad |w_{\delta,f}(u)|\leq\frac{\sqrt3}{2}\delta^2\|f\|_\infty.
\]
In the plane these hold for every compactly supported test.  A test
supported in a disk of physical radius $r$, with
$\|f\|_1\leq A$ and $\|f\|_\infty\leq Ar^{-2}$, therefore has
smearing weights at radius $a=r/\delta+O(1)$ with a constant multiple
of the same weight bound.  For any fixed physical geometry the
condition $a\geq1$ holds once the mesh is small.  Reflection in a
lattice symmetry takes each face to its reflected face, so this
interpolation is exactly equivariant under every site reflection used
below.  A zero-integral test gives a zero-total smearing at every mesh.

In particular, on a fixed compact subset of a bounded domain,
Proposition~\ref{prop:smearing} bounds a normalized radius-$r$ average
by $C_p(1+\sqrt{\log(1/r)})$, with fixed geometry and the mesh sent
to zero first.  Finite sums of such localized smearings obey the same
bounds by the triangle inequality.  If the radius is a fixed small
fraction of the distance to a bounded trace, the right side of
\eqref{moments:smearing-bound} is uniformly bounded, even as both
physical scales subsequently tend to zero.  Finally,
Proposition~\ref{prop:plane-moments} gives a uniform bound for each
fixed bounded family of zero-total smooth tests and their dilations.
These assertions use ordinary Euclidean area weights throughout.

The class hypothesis~\eqref{moments:class-central} has deliberately
not been asserted for any spin-pinned ensemble.  The parity-class
admissibility and central mass required in that setting will be proved
in Proposition~\ref{prop:pin-moments} before
Lemma~\ref{lem:midpoint-amplification} is applied there.
